% !TEX root = ../main.tex

\chapter{预备知识}

本章介绍后续算法设计、系统实现与部署验证所需的共同基础，重点包括导航任务形式、视觉目标导航、扩散模型与扩散策略、NoMaD 的关键机制，以及轮式平台与足式平台在部署上的差异。

\section{导航与视觉导航基础}

\subsection{导航任务的基本定义}

机器人导航的目标是在给定观察、任务目标和运动约束的条件下，使机器人从当前状态移动到目标区域。一个完整导航系统通常涵盖环境感知、状态表示、目标表达、局部决策与运动执行五个环节\cite{thrun2005probabilistic,anderson2018evaluation}。在输出形式上，高层导航策略可以输出底层控制量、局部路径点或局部轨迹序列；本文让高层模型输出局部航点轨迹，再由中层控制器完成执行映射。

如果把高层策略写成函数形式，则有
\begin{equation}
\pi : (\mathcal{O}_t, g) \mapsto \mathbf{A}_t ,
\end{equation}
其中 $\mathcal{O}_t$ 表示时刻 $t$ 的图像观察或观察序列，$g$ 表示任务目标，$\mathbf{A}_t$ 表示未来短时间内的局部轨迹。该表示天然适合与不同平台解耦，也是本文能够将高层视觉导航从轮式框架迁移到 Lite3 四足平台的重要原因。

观察空间方面，本文统一使用 RGB 图像作为输入，并保留最近 $K+1$ 帧上下文 $\mathcal{O}_t = \{I_{t-K}, \ldots, I_t\}$。每个航点$\mathbf{a}_i = (\Delta x_i, \Delta y_i)$ 表示相对于机器人当前位姿的二维增量位移（$x$ 前向、$y$ 左侧），策略输出因此不依赖全局定位。执行上采用闭环模式：每个高层周期生成长度 $H=8$ 的航点序列，固定取第 $k=2$ 个航点作为本次控制周期的局部目标位移，由中层 PD 控制器映射为单条 \texttt{Twist} 速度指令；下一周期再基于最新相机帧重新生成。这种``规划长 8 步、每周期只执行其中一个超前点''的策略与模型预测控制（MPC）一脉相承，扩散策略相当于把 MPC 的优化求解替换为条件扩散采样。

\subsection{视觉导航与目标导航}

视觉导航以图像作为主要环境输入，并基于视觉信息完成导航决策\cite{zhu2017targetdriven,anderson2018evaluation}。目标导航则进一步关注目标的给出形式，常见类别包括：
\begin{itemize}
    \item \textbf{point-goal 导航}\cite{ddppo2020}：目标以坐标点形式给出，要求机器人到达指定位置；
    \item \textbf{image-goal 导航}：目标以一张或多张图像给出，要求机器人到达图像所对应的位置；
    \item \textbf{object-goal 导航}：目标以物体类别给出（如``找到冰箱''），需要语义理解能力；
    \item \textbf{language-goal 导航}：目标以自然语言指令给出，是最灵活但也最具挑战的形式。
\end{itemize}
本文研究 image-goal 场景：相较于绝对坐标，图像目标更符合真实应用——用户更容易给出``去这个画面对应的位置''，而非统一坐标系下的精确终点；相较 object-goal 与 language-goal，图像目标无需语义理解模块、表达歧义更低，能在无 GPS 与无预建地图的场景中直接工作。

\subsection{拓扑导航与拓扑地图}

为了降低对精确几何地图的依赖，拓扑导航只保留若干代表性视觉节点及其连通关系\cite{savinov2018sptm,chaplot2020nts,shah2021ving}。本文系统中拓扑地图不是部署时由编码器在线生成的图结构，而是人工沿目标路径预先采集的一段\textbf{有序}视觉节点序列，文件名按 \texttt{0.png, 1.png, \ldots} 自然排序，节点编号同时编码``沿路径前进的方向''，因此拓扑地图在使用上更接近一条按空间顺序排好的导航队列。

设拓扑节点图像沿路径排列为 $\mathcal{G}=(I_g^{(0)}, \ldots, I_g^{(M-1)})$，导航任务给定起点附近索引 $0$ 与目标索引 $g\le M-1$。系统维护一个状态化的``当前定位''指针 $c_t$，距离预测头只在以 $c_t$ 为中心、半径 $r=4$ 的局部窗口
\begin{equation}
\mathcal{W}_t = \{j : \max(c_t-r,\,0) \le j \le \min(c_t+r+1,\,g)\}
\end{equation}
内对当前观察 $\mathcal{O}_t$ 与每个候选节点 $I_g^{(j)}$ 计算距离预测 $\hat d_j$，并据此更新定位：
\begin{equation}
c_{t+1} = \mathrm{clip}\bigl(\arg\min_{j\in\mathcal{W}_t} \hat d_j,\; c_t-1,\; c_t+\Delta_{\max}\bigr),
\end{equation}
其中 clip 操作确保定位指针每周期至多前进 $\Delta_{\max}=3$ 个节点、至多回退 1 个节点，从而强制定位整体沿拓扑地图单向推进。


\begin{figure}[htbp]
    \centering
    \includegraphics[width=0.9\linewidth]{prelim-nav-pipeline-placeholder.png}
    \caption{导航任务输入输出与拓扑地图关系示意图}
    \label{fig:prelim-nav-pipeline-placeholder}
\end{figure}

\section{扩散模型与扩散策略}

\subsection{扩散模型基本原理}

扩散模型通过``逐步加噪、逐步去噪''的方式学习复杂分布\cite{ddpm2020}。在前向过程里，数据样本 $x_0$ 被逐步扰动为高斯噪声：
\begin{equation}
q(x_t \mid x_{t-1}) = \mathcal{N}\!\left(\sqrt{1-\beta_t}x_{t-1}, \beta_t \mathbf{I}\right),
\end{equation}
其中 $\beta_t$ 为第 $t$ 步噪声调度系数。经过多步加噪后，样本会接近标准高斯分布。模型训练的目标是在反向过程中预测噪声，从而逐步恢复原始样本。常见训练目标写为
\begin{equation}
\mathcal{L}_{\text{diff}} = \mathbb{E}_{x_0,\epsilon,t}\left[\left\|\epsilon - \epsilon_\theta(x_t, t, c)\right\|_2^2\right],
\end{equation}
其中 $\epsilon$ 为真实噪声，$\epsilon_\theta$ 为模型预测噪声，$c$ 表示条件信息。

利用马尔可夫链的性质，可以推导出从 $x_0$ 到任意时刻 $x_t$ 的闭合形式。定义 $\alpha_t = 1 - \beta_t$，$\bar{\alpha}_t = \prod_{s=1}^{t} \alpha_s$，则有
\begin{equation}
q(x_t \mid x_0) = \mathcal{N}\!\left(\sqrt{\bar{\alpha}_t}\, x_0,\; (1 - \bar{\alpha}_t)\mathbf{I}\right).
\label{eq:closed-form-forward}
\end{equation}
这意味着可以直接从 $x_0$ 一步采样得到任意时刻的加噪样本 $x_t = \sqrt{\bar{\alpha}_t}\, x_0 + \sqrt{1 - \bar{\alpha}_t}\, \epsilon$，其中 $\epsilon \sim \mathcal{N}(0, \mathbf{I})$。这一闭合形式是训练过程高效进行的基础——训练时无需依次模拟整个前向链，只需随机采样时间步 $t$ 并直接计算 $x_t$ 即可。

本工作中显式将训练扩散步数设为 $T = 10$，这远低于通用图像扩散的 $T = 1000$ 量级，原因是本工作的扩散对象不是高维像素而是 16 维航点轨迹，低维分布对去噪步数的需求大幅降低，这一配置在轨迹生成质量与训练稳定性之间取得了良好的平衡。

从变分推断的角度来看，扩散模型的训练目标实际上是对数据对数似然的变分下界（Variational Lower Bound, VLB）的优化。具体而言，对数似然 $\log p_\theta(x_0)$ 的 VLB 可以分解为
\begin{equation}
\log p_\theta(x_0) \geq \mathbb{E}_q\!\left[-\log \frac{q(x_{1:T}\mid x_0)}{p_\theta(x_{0:T})}\right] = -\mathcal{L}_{\text{VLB}},
\end{equation}
其中 $\mathcal{L}_{\text{VLB}}$ 可进一步分解为逐步 KL 散度之和。Ho 等人\cite{ddpm2020}在实践中发现，直接使用简化的噪声预测损失 $\mathcal{L}_{\text{diff}}$ 训练效果更好，尽管它与严格的 VLB 目标并不完全等价。

在轨迹生成的具体应用中，$x_0$ 不再是图像像素，而是轨迹航点序列 $\mathbf{A} = [\mathbf{a}_1, \ldots, \mathbf{a}_H] \in \mathbb{R}^{H \times 2}$。相应地，扩散过程在轨迹空间中进行：前向过程逐步向轨迹添加噪声，反向过程从纯噪声逐步恢复出合理轨迹。由于轨迹空间的维度远低于图像空间，扩散模型在轨迹生成中的收敛速度和采样效率都显著优于图像生成，这也是扩散策略能够在实时控制场景中应用的重要原因。

这种建模方式在导航领域具有以下关键优势：
\begin{enumerate}
    \item \textbf{多模态建模能力}：扩散模型不强迫在多峰分布中只输出一个均值解，而是通过去噪过程逐步生成多种合理候选。在十字路口等存在多条合理路径的场景中，这种能力尤为重要；
    \item \textbf{样本质量可控}：通过调节去噪步数和引导强度，可以在生成质量与计算开销之间灵活权衡；
    \item \textbf{条件生成自然}：条件信息（如目标图像）可以自然地融入去噪过程，无需额外的条件注入机制。
\end{enumerate}
导航与控制任务中的未来轨迹本身就常常具有多模态特性，因此扩散建模天然适合此类问题。


\subsection{DDPM 与 DDIM}

DDPM 采用随机反向扩散过程生成样本，采样质量通常较好，但采样步数较多、推理开销较大\cite{ddpm2020}。DDIM 在保持训练目标不变的前提下，将反向过程改写为非马尔可夫且可确定性的更新形式，从而实现更少步数的快速采样\cite{ddim2020}。当取确定性采样形式时，其一步更新可写为
\begin{equation}
x_{t-1} = \sqrt{\bar\alpha_{t-1}}\hat x_0 + \sqrt{1-\bar\alpha_{t-1}}\,\epsilon_\theta(x_t,t,c),
\end{equation}
其中 $\hat x_0$ 由当前样本 $x_t$ 与预测噪声反推得到，$\bar\alpha_t$ 为累计噪声调度系数。

更具体地，DDIM 的关键创新在于引入了一族非马尔可夫前向过程。Song 等人\cite{ddim2020}观察到，DDPM 的训练目标仅依赖于边际分布 $q(x_t \mid x_0)$ 而非联合分布 $q(x_{1:T} \mid x_0)$。因此可以构造一族具有相同边际分布但不同联合分布的前向过程，其反向过程可以写为
\begin{equation}
q_\sigma(x_{t-1} \mid x_t, x_0) = \mathcal{N}\!\left(\sqrt{\bar\alpha_{t-1}}\, x_0 + \sqrt{1-\bar\alpha_{t-1}-\sigma_t^2}\cdot\frac{x_t - \sqrt{\bar\alpha_t}\, x_0}{\sqrt{1-\bar\alpha_t}},\; \sigma_t^2 \mathbf{I}\right),
\end{equation}
其中 $\sigma_t$ 是一个可调参数，控制反向过程中的随机性。定义 $\eta$ 参数使得
\begin{equation}
\sigma_t =
\eta
\sqrt{\frac{1-\bar\alpha_{t-1}}{1-\bar\alpha_t}}
\sqrt{1-\frac{\bar\alpha_t}{\bar\alpha_{t-1}}},
\end{equation}
则 $\eta = 1$ 时恢复为 DDPM 的随机采样，$\eta = 0$ 时得到完全确定性的 DDIM 采样。$\eta$ 的中间值则在两者之间插值，提供了在采样多样性与确定性之间的连续调节能力。

在导航轨迹生成的实际应用中，$\eta$ 的选择涉及探索与确定性之间的权衡。$\eta = 0$（完全确定性）意味着给定相同的初始噪声和条件，每次都会生成完全相同的轨迹，这在需要可复现行为的场景中有优势，但丧失了多样性。$\eta > 0$ 则引入随机性，使得每次采样可能产生不同的轨迹候选，在本文实验中，$\eta$ 被设置为 $\eta \in [0.3, 0.5]$），以确保候选轨迹之间有足够的多样性。

\begin{table}[htbp]
    \centering
    \caption{DDPM 与 DDIM 的主要性质对比}
    \label{tab:ddpm-ddim-compare}
    \begin{tabular}{lll}
        \toprule
        性质 & DDPM & DDIM \\
        \midrule
        反向过程类型 & 随机（马尔可夫） & 确定性或半随机（非马尔可夫） \\
        采样步数 & 通常需要完整 $T$ 步 & 支持子序列跳步采样 \\
        $\eta$ 参数 & 固定为 1 & 可调，0 为确定性 \\
        采样多样性 & 较高（每次结果不同） & $\eta=0$ 时完全确定 \\
        训练目标 & 噪声预测 $\mathcal{L}_{\text{diff}}$ & 与 DDPM 相同，无需重新训练 \\
        推理速度 & 较慢 & 可显著加速 \\
        \bottomrule
    \end{tabular}
\end{table}

本文后续所说的``DDIM 推理加速''，就是在不重新训练网络的前提下，用更少的去噪步数近似原始 DDPM 采样。

在实际使用中，DDIM 的跳步采样通过选择原始时间步序列 $\{1, 2, \ldots, T\}$ 的一个子序列 $\{\tau_1, \tau_2, \ldots, \tau_S\}$（其中 $S \le T$）来实现加速。NoMaD 训练时已经把 $T$ 压到 $10$，因此本文 DDIM 实验的有效压缩区间是从 $S = 10$ 进一步减到 $S = 5/3/2/1$，对应每步覆盖 $1$ 到 $5$ 个原始时间步。子序列的选择方式（均匀间隔、二次间隔等）也会影响采样质量。在本文的导航场景中，由于轨迹空间的维度远低于图像空间，扩散模型对跳步采样的容忍度通常较高，这为在边缘设备上实现实时推理提供了可能。

\subsection{Classifier-Free Guidance}

Classifier-Free Guidance（CFG）最初用于条件生成任务中的引导增强\cite{cfg2022}。其思想是在训练阶段同时学习有条件与无条件预测，在推理阶段通过两者线性组合增强条件约束。若条件分支噪声预测为 $\epsilon_{\text{cond}}$，无条件分支为 $\epsilon_{\text{uncond}}$，则引导后的噪声预测写为
\begin{equation}
\epsilon_{\text{cfg}} = (1+w)\epsilon_{\text{cond}} - w\epsilon_{\text{uncond}},
\end{equation}
其中 $w$ 为 guidance scale。当 $w=0$ 时退化为普通条件推理；当 $w>0$ 时，模型会更强地朝向条件约束；当 $w$ 过大时，也可能带来分布偏移与不稳定。

CFG 之所以能在推理时实现无分类器引导，关键在于训练阶段的条件随机丢弃（condition dropout）机制。具体而言，在训练过程中以概率 $p_{\text{uncond}}$（通常取 $0.1$ 至 $0.2$）将条件信息 $c$ 替换为空条件 $\varnothing$，使得同一模型 $\epsilon_\theta$ 同时学习条件预测 $\epsilon_\theta(x_t, t, c)$ 与无条件预测 $\epsilon_\theta(x_t, t, \varnothing)$。在推理阶段，对同一输入分别执行条件和无条件前向传播，
再通过线性组合实现引导，无需额外训练分类器网络。


在本文中，CFG 的作用是增强轨迹对目标条件的牵引能力，即让生成的局部航点更明确地朝向目标图像或目标节点图像对应的方向。

本文中 CFG 的推理过程可概括如下：
\begin{enumerate}
    \item 在每个去噪步骤中，分别以目标 token 和掩码 token 作为条件，执行两次前向传播；
    \item 获得条件噪声预测 $\epsilon_{\text{cond}}$ 和无条件噪声预测 $\epsilon_{\text{uncond}}$；
    \item 按照 CFG 公式计算引导后的噪声预测 $\epsilon_{\text{cfg}}$；
    \item 使用 $\epsilon_{\text{cfg}}$ 执行一步去噪更新。
\end{enumerate}
需要注意的是，在实际实现中，CFG 的无条件分支并非完全无信息——它仍然接收历史观察 token，只是目标 token 被替换为掩码 token。因此 CFG 在这里增强的是``目标方向性''，而非``整体行为合理性''。这一区别使得即使在较大的 $w$ 值下，生成的轨迹仍然保持基本的可行性（因为观察条件始终存在），只是目标趋向性被放大。

\subsection{扩散策略与 Test-Time Scaling}

扩散策略将扩散模型从图像生成推广到动作序列或轨迹序列生成\cite{diffuser2022,chi2025diffusion}。在具身智能与机器人控制领域，Diffuser 和 Diffusion Policy 证明了扩散建模能够有效处理多步决策、长时序动作和多模态行为分布，因此被视为扩散模型在具身控制中的代表性应用。

具体而言，Diffuser\cite{diffuser2022} 将整个轨迹（包括状态和动作序列）作为扩散对象，在推理时通过引导函数注入奖励信息来生成高回报轨迹。Diffusion Policy\cite{chi2025diffusion} 则只对动作序列进行扩散，将状态观察作为条件信息，更适合视觉观察驱动的控制任务。NoMaD 的扩散策略头采用了类似 Diffusion Policy 的设计——仅对航点轨迹进行扩散，视觉特征作为条件信息注入。这种设计使得扩散过程仅需在低维轨迹空间中进行，显著降低了计算开销。

在此基础上，近年来出现了测试时扩展（Test-Time Scaling, TTS）的思路\cite{ttsdiffusion2025,ttssearch2025}。其核心思想不是修改训练过程，而是在推理阶段使用额外计算预算生成多条候选轨迹，再通过打分函数进行选择。若生成 $N$ 条候选轨迹 $\{\mathbf{A}^{(i)}\}_{i=1}^N$，评分函数为 $s(\cdot)$，则最优轨迹可表示为
\begin{equation}
\mathbf{A}^{\ast} = \arg\max_{i\in\{1,\ldots,N\}} s\!\left(\mathbf{A}^{(i)} \mid \mathcal{O}_t, I_g\right).
\end{equation}

评分函数的设计是 TTS 方法的核心。常见的评分维度包括：（1）目标一致性评分，衡量候选轨迹终点与目标方向的对齐程度，可以利用距离预测头对轨迹终点的距离估计作为代理指标；（2）轨迹平滑性评分，衡量相邻航点之间的方向变化和速度变化是否连续，过大的方向突变通常意味着不可执行或执行风险较高的轨迹；（3）前向进展评分，衡量轨迹整体是否朝目标方向前进，避免选择在原地打转或后退的候选。在实际实现中，这些评分维度通常以加权和的形式组合：
\begin{equation}
s(\mathbf{A}^{(i)}) = \lambda_{\text{goal}} \cdot s_{\text{goal}}(\mathbf{A}^{(i)}) + \lambda_{\text{smooth}} \cdot s_{\text{smooth}}(\mathbf{A}^{(i)}) + \lambda_{\text{prog}} \cdot s_{\text{prog}}(\mathbf{A}^{(i)}),
\end{equation}
其中 $\lambda_{\text{goal}}$、$\lambda_{\text{smooth}}$ 和 $\lambda_{\text{prog}}$ 为各维度的权重系数。具体地，平滑性评分可以定义为相邻航点方向变化角的负值之和：
\begin{equation}
s_{\text{smooth}}(\mathbf{A}) = -\sum_{i=2}^{H-1} \left|\angle(\mathbf{a}_{i+1} - \mathbf{a}_i) - \angle(\mathbf{a}_i - \mathbf{a}_{i-1})\right|,
\end{equation}
其中 $\angle(\cdot)$ 表示二维向量的方向角。前向进展评分则可以定义为轨迹终点在目标方向上的投影距离：
\begin{equation}
s_{\text{prog}}(\mathbf{A}) = \mathbf{a}_H \cdot \hat{\mathbf{d}}_{\text{goal}},
\end{equation}
其中 $\hat{\mathbf{d}}_{\text{goal}}$ 为指向目标的单位方向向量。

TTS 的工作方式与蒙特卡洛树搜索（Monte Carlo Tree Search, MCTS）存在概念上的相似性。MCTS 通过大量随机模拟评估不同决策路径的期望回报，TTS 则通过大量采样评估不同生成轨迹的综合质量。两者的核心共性在于``用额外计算换取更优决策''。区别在于 MCTS 通常涉及多步展开和回溯，而 TTS 在扩散策略中主要针对单次生成的多条候选进行比较选择，其计算开销的扩展方式更为简单和可控。

从理论角度来看，TTS 的有效性建立在一个关键假设之上：扩散模型生成的候选集合中，高质量轨迹的存在概率随候选数量 $N$ 的增加而增大。设单次采样得到可接受轨迹的概率为 $p$，则 $N$ 次独立采样后至少存在一条可接受轨迹的概率为 $1 - (1-p)^N$。当 $p$ 较小时（对应困难场景），需要较大的 $N$ 才能保证成功率；当 $p$ 接近 1 时（对应简单场景），$N=1$ 即可，此时 TTS 退化为普通单次采样。这种``按需扩展''的
特性使得 TTS 特别适合于计算预算可变的部署场景。

需要强调的是，TTS 与 DDIM、CFG 作用位置不同。DDIM 决定``怎么采样''，CFG 决定``采样时如何增强条件''，而 TTS 决定``采样出多条候选后如何选择''。因此三者虽然都属于推理时优化，但不是简单相加关系，而是具有先后层次的组合关系。

\subsection{扩散模型、ViT 与导航建模的关系}

Transformer 在视觉建模中的广泛应用，使得扩散模型与视觉编码方法的结合成为趋势。ViT 将图像切分为 patch token 并以自注意力机制进行建模\cite{vit2020}；DiT 则将扩散过程进一步与 Transformer 架构结合，说明扩散模型可以在强表征模型之上获得更高质量生成能力\cite{dit2022}。在导航领域，ViNT 使用 Transformer 处理历史观察与目标图像\cite{vint2023}，NoMaD 则在此基础上引入动作扩散策略\cite{nomad2023}。因此，本文讨论的视觉编码器升级、扩散采样优化和系统部署，其实都位于这一技术演化脉络之中。

具体而言，ViT 的核心操作是将输入图像划分为固定大小的 patch（如 $16 \times 16$ 像素），将每个 patch 线性映射为一个 token 向量，再加上位置编码后送入多层 Transformer 编码器。对于输入尺寸为 $H \times W$ 的图像和 patch 大小 $P$，token 序列长度为 $N = HW/P^2$。

\begin{figure}[htbp]
    \centering
    \includegraphics[width=0.84\linewidth]{prelim-diffusion-placeholder.png}
    \caption{扩散模型、扩散策略、ViT 编码器与 NoMaD 系统之间的技术演化关系。
    从左到右依次为基础扩散理论（DDPM/DDIM）、条件引导方法（CFG）、
    扩散策略（Diffusion Policy）以及 NoMaD 的架构}
    \label{fig:prelim-diffusion-placeholder}
\end{figure}

\section{NoMaD 机制}

\subsection{NoMaD 的基本结构}

NoMaD 由视觉编码器、距离预测头和扩散策略头三部分构成\cite{nomad2023}。视觉编码器负责把历史观察与目标图像编码为统一条件特征；距离预测头估计当前观察与目标之间的相对距离
（以离散时间步数为单位，表示从当前位置到达目标大约需要多少个导航步骤，而非物理距离）；扩散策略头则在该条件特征上生成未来局部航点轨迹。若记局部轨迹为
\begin{equation}
\mathbf{A} = [\mathbf{a}_1,\mathbf{a}_2,\ldots,\mathbf{a}_H], \qquad \mathbf{a}_i=(\Delta x_i,\Delta y_i),
\end{equation}
则 NoMaD 学习的是条件分布
\begin{equation}
p(\mathbf{A}\mid \mathcal{O}_t, I_g).
\end{equation}

在 NoMaD 的视觉编码器部分，其底层基于 ViNT 的 Transformer 架构。具体而言，每帧观察图像首先通过卷积神经网络（Convolutional Neural Network, CNN）或视觉 Transformer（Vision Transformer, ViT）前端提取特征图，再经过全局池化与线性投影映射为统一的 token 表示。若历史上下文长度为 $K$，则模型将当前帧及其之前的 $K$ 帧图像编码为观察 token；目标图像与当前帧共同编码为目标 token。这些 token 被拼接为一个统一的序列，送入 Transformer 编码器进行联合注意力计算。形式化地，令观察 token 序列为 $\{\mathbf{z}_{\text{obs}}^{(k)}\}_{k=0}^{K}$，目标 token 为 $\mathbf{z}_{\text{goal}}$，则 Transformer 编码器的输入为
\begin{equation}
\mathbf{Z}_{\text{input}} = [\mathbf{z}_{\text{obs}}^{(0)}; \mathbf{z}_{\text{obs}}^{(1)}; \ldots; \mathbf{z}_{\text{obs}}^{(K)}; \mathbf{z}_{\text{goal}}],
\end{equation}
共 $K+2$ 个 token。经过多层自注意力后，输出的 token 表征即为共享条件空间中的统一条件特征。

在 ViNT\cite{vint2023} 的原始设计中，Transformer\cite{transformer2017} 编码器采用多层标准自注意力结构。每帧图像经过 EfficientNet-B0\cite{efficientnet2019} 前端后形成图像特征，再由线性投影映射为 Transformer 可处理的 token。为了区分不同时间步的观察，模型为各帧 token 添加帧级位置编码。Transformer 的输出 token 经过池化后产生固定维度的条件向量，该向量同时被距离预测头和扩散策略头使用，因此视觉编码器的表征质量会同时影响目标距离估计与局部轨迹生成。

\subsection{距离预测头与视觉编码器的直接耦合}

距离预测头与扩散策略头共享同一条视觉条件空间，这意味着视觉编码器的改动会产生连锁效应。假设原始编码器输出的条件特征维度为 $d_{\text{orig}}$，新编码器输出维度为 $d_{\text{new}}$，则即使通过线性投影层将 $d_{\text{new}}$ 映射到 $d_{\text{orig}}$，投影后的特征分布也与原始特征不同。距离预测头是一个在原始特征分布上训练的回归网络，输入分布的偏移会直接导致距离预测不准确。因此，本文在编码器升级实验中采用端到端微调策略，同时更新编码器适配层、距离预测头和扩散策略头的参数。这种端到端微调策略虽然计算成本较高，但确保了三个组件在新的特征空间中重新对齐，避免了因局部更新导致的性能退化。

\subsection{goal mask 与导航/探索统一建模}

NoMaD 的另一个关键机制是 goal mask。训练时，模型以一定概率遮蔽目标 token，使其同时经历``有目标条件''和``无显式目标条件''两种情况\cite{nomad2023}。这并不意味着模型在无目标条件时进行随机游走，而是意味着它会学习在当前视觉上下文中生成连续、可执行且符合数据轨迹分布的探索性行为。正因为如此，NoMaD 才能在部署时同时支持导航模式与探索模式。

goal mask 的设计还有一个重要的技术意义：它使得模型在训练时自然地学习了
``当前观察本身蕴含的行为先验''。即使没有目标指引，训练数据中的轨迹通常
也遵循特定的行为模式（如沿道路行驶、避开障碍物），goal mask 使模型能够
将这些模式内化为无条件生成能力。

在数学表述上，goal mask 的训练过程可以描述如下。设 goal mask 概率为 $p_{\text{mask}}$，对于每个训练样本，以概率 $p_{\text{mask}}$ 将目标 token $\mathbf{z}_{\text{goal}}$ 替换为可学习的掩码 token $\mathbf{z}_{\text{mask}}$。因此模型在训练过程中优化的是混合目标：
\begin{equation}
\mathcal{L} = p_{\text{mask}} \cdot \mathcal{L}_{\text{diff}}(\epsilon_\theta(x_t, t, \mathbf{z}_{\text{obs}}, \mathbf{z}_{\text{mask}})) + (1 - p_{\text{mask}}) \cdot \mathcal{L}_{\text{diff}}(\epsilon_\theta(x_t, t, \mathbf{z}_{\text{obs}}, \mathbf{z}_{\text{goal}})).
\end{equation}
在本文的配置中，$p_{\text{mask}}$ 设置为 $0.5$ ，确保两种模式都获得充分训练。值得注意的是，goal mask 机制与 CFG 中的条件随机丢弃（condition dropout）在形式上高度一致——两者都是在训练时以一定概率移除条件信息。因此 goal mask 实际上内在地为 CFG 推理提供了基础：在推理时，将条件分支（提供目标 token）与无条件分支（使用掩码 token）的噪声预测进行 CFG 组合，即可实现对目标方向的引导增强。这种设计使得探索能力与导航增强能力共享同一训练框架。

在部署阶段，探索模式与导航模式的切换仅需改变推理时的条件输入：提供真实目标 token 时进入导航模式，提供掩码 token 时进入探索模式。两种模式之间的切换不需要加载不同的模型权重或改变网络结构，这大大简化了部署系统的设计。此外，探索模式下生成的轨迹并非随机——它们反映了训练数据中的轨迹分布特征，通常表现为沿走廊前进、绕过障碍物等符合环境结构的行为模式。

\section{轮式平台、足式平台与 Lite3}

\subsection{轮式与足式平台的部署差异}

轮式平台姿态变化小、视觉输入稳定，与速度控制器的接口也较简单；足式平台则具有更强的地形适应性，但会带来更明显的机体起伏、更严格的实时性约束和更复杂的控制接口\cite{hwangbo2019anymal,lee2020learning}。这意味着，同一高层轨迹若能在轮式平台上顺利执行，并不代表它能直接在四足平台上稳定工作。

具体而言，足式平台的步态运动会在每个步态周期内引入机体俯仰和侧倾变化，这种扰动会导致相机的视场发生抖动，使得连续帧之间的视觉特征产生非平移性变化。对于依赖时序视觉输入的导航策略，步态引起的视觉抖动可能被误解读为环境变化或自身运动，从而产生错误的轨迹预测。此外，足式平台的非完整运动学约束也与轮式平台不同。轮式机器人通常可以精确执行差速运动学模型所描述的运动，而四足机器人的实际运动轨迹受步态规划器和地形交互的影响，与指令速度之间存在非线性偏差。这种偏差意味着即使高层策略输出了精确的速度指令，机器人的实际位移也可能与预期存在差异，进一步强调了闭环执行和频繁重规划的必要性。
 

\subsection{云深处 Lite3 平台与关键参数}

本文采用的主要足式平台为云深处 Lite3。结合本文系统设计与部署需求，
其在本文中的系统定位是``由高层视觉导航策略驱动的四足运动系统''，而不是端到端关节控制平台。对应控制链路为
\begin{equation}
\text{视觉观察} \rightarrow \text{高层 NoMaD 推理} \rightarrow \text{中层控制映射} \rightarrow \text{底层运动执行}.
\end{equation}

与本文实现直接相关的关键参数如表~\ref{tab:lite3-params} 所示。
\begin{table}[htbp]
    \centering
    \caption{本文中与 Lite3 相关的关键平台参数}
    \label{tab:lite3-params}
    \begin{tabular}{lll}
        \toprule
        参数类别 & 参数项 & 数值或说明 \\
        \midrule
        平台类型 & 机器人形态 & 四足平台 \\
        运动结构 & 关节数量 & 12 关节执行系统 \\
        速度接口 & 前后速度范围 & $[-1.0,1.0]\,\text{m/s}$ \\
        速度接口 & 左右速度范围 & $[-0.5,0.5]\,\text{m/s}$ \\
        速度接口 & 偏航角速度范围 & $[-1.5,1.5]\,\text{rad/s}$ \\
        部署限幅 & 本文最大线速度 & $0.12\,\text{m/s}$ \\
        部署限幅 & 本文最大角速度 & $0.35\,\text{rad/s}$ \\
        高层闭环 & 仿真 NoMaD 周期 & $4\,\text{Hz}$ \\
        通信安全 & 速度超时保护 & $250\,\text{ms}$ 自动停止 \\
        \bottomrule
    \end{tabular}
\end{table}


\begin{figure}[htbp]
    \centering
    \includegraphics[width=0.5\linewidth]{prelim-lite3-placeholder.png}
    \caption{Lite3 平台与控制链路示意图}
    \label{fig:prelim-lite3-placeholder}
\end{figure}

在通信协议层面，本文系统与 Lite3 之间通过 UDP 协议进行数据交换。桥接层将高层三维速度参考转换为前后、左右和偏航三个速度通道，并持续维护心跳与速度刷新。本文设置心跳频率为 $5\,\text{Hz}$，速度刷新频率为 $25\,\text{Hz}$，满足 Lite3 通讯文档中速度轴命令至少 $20\,\text{Hz}$ 刷新的要求。协议文档给出的速度命令超时约束可以理解为底层接口的安全边界，但高层 NoMaD 推理本身并不需要在每个 $250\,\text{ms}$ 内都完整完成一次；桥接层会在两次高层推理之间持续重复最近一次限幅后的速度参考。

关于观察频率与控制频率的关系，需要区分三个层次的频率，这种多层频率架构可以用以下层级关系概括：
\begin{equation}
f_{\text{joint}} \gg f_{\text{velocity}} \gg f_{\text{vision}},
\end{equation}
其中 $f_{\text{joint}}$ 表示底层关节级控制频率，$f_{\text{velocity}}$ 表示速度命令刷新频率，当前真机桥接实现中 $f_{\text{velocity}}=25\,\text{Hz}$，$f_{\text{vision}}$ 表示高层视觉导航策略的更新频率，在 MuJoCo 仿真中为 $4\,\text{Hz}$，在 Orin 真机上则由实测推理耗时决定。这种频率分离设计的好处是每一层只需关注自身时间尺度上的决策，坏处是层间的响应延迟会累积。  
